\naslov{Premočrtno gibanje}

\begin{definicija}
  Gibanje je \pojem{premočrtno}, če ima pospešek konstantno smer.
\end{definicija}

Primer takega gibanja je poševni met.
Opazimo, da lahko vedno izberemo koordinatni sistem, v katerem tir poti leži na
premici:
Če je $\vec{a} = a \vec{e}$, kjer je $\vec{e}$ konstanten vektor, velja
\[
  \vec{v} = \vec{e} \int_{t_0}^t a dt + \vec{v}_0.
\]
Izberemo lahko sistem, kjer je $\vec{v}_0$ enak $\vec{0}$, in bo torej $\vec{v}$
vzporeden $\vec{e}$.

Če gibanje poteka pod vplivom konzervativne sile, lahko zapišemo potencial $U$,
in velja izrek o energiji
\[
  \pol m \dot{x}^2 + U(x) = E_0.
\]
Od tod izpeljemo
\[
  \dot{x} = \pm \sqrt{\frac{2}{m} \left( E_0 - U(x) \right) }.
\]
Enačbo z ločljivimi spremenljivkami tedaj integriramo in dobimo
\[
  \pm \int_{x_0}^{x} \frac{dx}{\sqrt{ \frac{2}{m} \left( E_0 - U(x) \right) } }
  = \int_{t_0}^t dt = t - t_0.
\]
Dobimo funkcijo $t = t(x)$.
Če se na poti ne ustavimo, po izreku o inverzni preslikavi obstaja funkcija $x
= x(t)$.
Pravimo, da je premočrtno gibanje \pojem{integrabilno}.

Če v kvalitativni analizi ugotovimo, da je neko gibanje periodično med točkama
$a$ in $b$, lahko periodo izračunamo kot
\begin{align*}
  T &= \int_{x_0}^b \frac{dx}{\sqrt{ \frac{2}{m} \left( E_0 - U(x) \right) }}
  - \int_a^b \frac{dx}{\sqrt{ \frac{2}{m} \left( E_0 - U(x) \right) }}
  + \int_a^{x_0} \frac{dx}{\sqrt{ \frac{2}{m} \left( E_0 - U(x) \right) }} \\
  &= \sqrt{2 m} \int_a^b \frac{dx}{ \sqrt{ E_0 - U(x) } }.
\end{align*}
Ker je $E = U(x)$ v krajiščih, je to posplošen integral.
Situacija v obeh krajiščih je simetrična, torej preverimo le za levo krajišče,
da integral res konvergira.
V prvem koraku razvijemo preslikavo $U$ v Taylorjev polinom prve stopnje v točki
$a$, kjer se pojavi vrednost odvoda $U$ v neki točki $\xi$ blizu $a$.
Ker je odvod zvezen, obstaja tak $\delta > 0$, da za $\xi \in [a, a + \delta)$
velja $2 \partial_x U(a) < \partial_x U(\xi) < \pol \partial_x U(a)$, torej
\[
  \int_a^{a + \delta} \frac{dx}{\sqrt{E_0 - U(x)}}
  = \int_a^{a + \delta} \frac{dx}{\sqrt{ - \partial_x U(\xi) (x-a) }}
  \le \int_a^{a + \delta} \inv{\sqrt{ -\pol \partial_x U(a) }} \inv{\sqrt{x -
	  a}} dx
  < \infty.
\]

\vprasanje{Izpelji izraz za periodo premočrtnega potencialnega gibanja.}

\begin{lema}
  Za $a < b$ velja
  \[
	\int_a^b \frac{dx}{\sqrt{(b-x) (x-a)}} = \pi.
  \]
\end{lema}

\begin{proof}
  Uvedemo novo spremenljivko $x = \pol (a+b) + \pol (b-a) z$, s čimer se
  integral spremeni v
  \[
	\int_{-1}^1 \frac{dz}{\sqrt{(1-z)(1+z)}}
	= \pi.
  \]
\end{proof}

\begin{primer}
  Oglejmo si harmonični oscilator, ki deluje pod potencialom $U = \pol k x^2$.
  Tedaj velja $F = - \partial_x U = -kx$, torej $m\ddot{x} = -kx$.
  Rešitev tega sistema je $x = A \cos \omega t + B \sin \omega t$ za
  $\omega = \sqrt{k/m}$.
  Iz tega lahko kar direktno preberemo $T = 2 \pi / \omega$.
  Posebnost harmoničnega oscilatorja je, da je $T$ neodvisen od $E_0$.
  Takemu gibanju pravimo \pojem{izohronično}, harmonični potencial je edini
  primer izohroničnega potenciala, ki je simetričen glede na svoj minimum.
\end{primer}

Če ima potencial lokalni minimum v $x_0$, lahko za določanje potenciala
uporabimo harmonično aproksimacijo.
Zapišemo
\[
  \hat{U}(x) = U(x_0) + \partial_x U(x_0) (x-x_0) + \pol \partial_{x^2}
  U(x_0) (x - x_0)^2,
\]
kar je harmonični potencial s periodo
\[
  T = 2 \pi \sqrt{\frac{m}{\partial_{x^2} U(x_0)}}.
\]
Ta aproksimacija je dobra, če velja $E_0 - U(x) \ll 1$.

\vprasanje{Izpelji harmonično aproksimacijo.}

Druga vrsta aproksimacije, ki jo lahko uporabimo, je \pojem{libracijska}.
Računamo
\[
  t = \sgn \dot{x} \int_{x_0}^x \frac{dx}{\sqrt{ \frac{2}{m} (E_0 - U(x)) }}
  = \sqrt{\frac{m}{2}} \sgn \dot{x}
  \int_{\theta_0}^{\theta} \frac{\pol (b-a) (-\sin \theta) d\theta}{\pol (b-a)
	\sqrt{\chi(\theta)} \sqrt{1 - \cos^2 \theta}}
\]
za substitucijo $x = \pol (a+b) + \pol (b-a) \cos\theta$ in $E_0 - U(x) =
(x-a)(b-x) \chi(x)$.
Pri tem smo si $x$ predstavljali kot kosinus kota v krožnici, ki poteka skozi
točki $a$ in $b$ in ima središče na njuni zveznici.
Če računamo dalje, dobimo
\[
  t = \sqrt{\frac{m}{2}} \int_{\theta_0}^{\theta} \inv{\sqrt{\chi(\theta)}}
  d\theta.
\]
Če želimo dobiti periodo gibanja, bo $\theta$ tekel od $0$ do $\pi$.
\[
  T = 2 \sqrt{\frac{m}{2}} \int_0^\pi \frac{d\theta}{\sqrt{\chi(\theta)}}.
\]
Ta integral aproksimiramo s trapezno formulo, ki je natančna v primeru, da je
funkcija v integralu afina.
Rešitev je tedaj
\[
  T \doteq \pi \sqrt{\frac{m}{2}} \left( \inv{\sqrt{\chi(a)}} +
	\inv{\sqrt{\chi(b)}} \right).
\]

\vprasanje{Izpelji libracijsko aproksimacijo.}

\begin{trditev}
  Za premočrtno potencialno periodično gibanje velja $\frac{dS}{dE_0} =
  \frac{T}{m}$ za ploščino $S$ faznega diagrama.
\end{trditev}

\begin{proof}
  Velja $S = 2 \sqrt{\frac{m}{2}} \int_a^b \sqrt{E_0 - U} dx$, torej
  \[
	\frac{dS}{dE_0} = 2 \sqrt{\frac{m}{2}} \left(
	  b' \sqrt{E_0 - U(b)} - a' \sqrt{E_0 - U(a)} + \int_a^b
	  \frac{dx}{2\sqrt{E_0 - U}}
	\right).
  \]
  Ker je $U(a) = U(b) = E_0$, sta prva dva člena v oklepaju enaka $0$, torej
  \[
	\frac{dS}{dE_0} = \sqrt{\frac{m}{2}} \int_a^b \frac{dx}{\sqrt{E_0-U}}
	= \frac{T}{m}.
  \]
\end{proof}

\vprasanje{Kako se pri nihanju ploščina faznega portreta spreminja z energijskim
  nivojem $E_0$?}

\naslov{Gibanje po krivulji}

Dana je krivulja $\vec{r} = \vec{r}(s(t))$, kjer je $s$ naravni parameter.
Če s $\mb{P}$ označimo trenutno lokacijo, velja
\[
  \vec{v} = \frac{d\mb{P}}{dt} = \frac{d\mb{P}}{ds} \frac{ds}{dt} = \vec{e}_t
  \dot{s},
\]
kjer je $\vec{e}_t$ enotski vektor, tangenten na krivuljo, in
\[
  \vec{a} = \ddot{s} \vec{e}_t + \dot{s} \frac{d \vec{e}_t}{dt}
  = \ddot{s} \vec{e}_t + \dot{s}^2 \kappa \vec{e}_n.
\]
V enačbi $\kappa$ predstavlja ukrivljenost, $\vec{e}_n$ pa normalo na krivuljo.
Drugi Newtonov zakon poleg rezultante vseh sil vsebuje tudi silo vezi
$\vec{S}$.
Razpisan v smereh krivuljnega koordinatnega sistema ima obliko
\begin{gather*}
  m \ddot{s} = \vec{F} \cdot \vec{e}_t + \vec{S} \cdot \vec{e}_t \\
  m \kappa \dot{s}^2 = \vec{F} \cdot \vec{e}_n + \vec{S} \cdot \vec{e}_n \\
  0 = \vec{F} \cdot \vec{e}_b + \vec{S} \cdot \vec{e}_b
\end{gather*}
Tu imamo štiri neznanke ($s, \vec{S}$) ter tri enačbe, torej potrebujemo še
dodatno konstitutivno relacijo za silo vezi.
Če se omejimo na gladke krivulje (take, kjer ni trenja), dobimo dodatno enačbo
\[
  \vec{S} \cdot \vec{e}_t = 0.
\]
Delo take sile vezi je enako $0$.
Če je $\vec{F}$ konzervativna sila, $\vec{F} = - \grad U$, dobimo
\[
  m \ddot{s} = - \frac{dU}{ds},
\]
iz česar lahko izpeljemo energijsko enačbo
\[
  \pol m \dot{s}^2 + U(s) = E_0.
\]
Torej je gibanje po gladki krivulji pod vplivom konzervativne sile reducibilno
na premočrtno gibanje v ločni dolžini.

\vprasanje{Na kaj se reducira gibanje po gladki krivulji pod vplivom
  konzervativne sile? Izpelji.}

\begin{vo}{Obravnavaj matematično nihalo kot gibanje po gladki krožnici.}
  Če je $l$ polmer krožnice, velja $s = l \theta$.
  Na točko poleg sile vezi deluje tudi teža, ki ima potencial
  \[
	U = - m \vec{g} \vec{r} = - m g l \cos \frac{s}{l}.
  \]
  Za dovolj majhen $E_0$ je gibanje periodično, in velja
  \[
	T = \sqrt{2m} \int_{-s_0}^{s_0} \frac{ds}{\sqrt{E_0 + mgl \cos
		\frac{s}{l}}}.
  \]
  Če uporabimo substitucijo $s = l \theta$ in začetno energijo zapišemo z
  začetnim odklonom, dobimo
  \[
	T = \sqrt{2m} l \int_{-\theta_0}^{\theta_0} \frac{d\theta}{mgl(\cos \theta -
	  \cos \theta_0)}.
  \]
  Na tej točki upoštevamo, da je funkcija soda, in uporabimo $\cos \theta = 1 -
  2 \sin^2 \frac{\theta}{2}$;
  \[
	T = 2 \sqrt{\frac{l}{g}} \int_0^{\theta_0} \frac{d\theta}{\sqrt{ \sin^2
		\frac{\theta_0}{2} - \sin^2 \frac{\theta}{2} }},
  \]
  kar se s substitucijo $\sin \frac{\theta}{2} = u \sin{\theta_0}{2}$ končno
  predela na
  \[
	 T = 4 \sqrt{\frac{l}{g}} \int_0^1 \frac{du}{\sqrt{ (1-u^2) (1 - u^2 \sin^2
		\frac{\theta_0}{2}) }}.
  \]
  To je eliptični integral, odgovor je
  \[
	T = 4 \sqrt{\frac{l}{g}} K\left( \sin^2 \frac{\theta_0}{2} \right)
  \]
\end{vo}

% LocalWords:  reducibilno
